%%%%%%%%%%%%%%%%%%%%%%%%%%%%%%%%%%%%%%%%%%%%%%%%%%%%%%%%%%%%%%%%%%%%%%%%%%%%%%%
\section{Published signals with known fates}
\label{sec:bench}

\subsection{Construction}

We compiled \nmBenchN\ periodic signals once claimed as planets around HARPS-RVBank
stars with at least 20 nights, and labelled each by the later peer-reviewed verdict
(Table~\ref{tab:bench}): \nmBenchUI\ upheld by independent evidence (transits, another
spectrograph or an independent data set), \nmBenchUH\ accepted from HARPS data and not
contested, \nmBenchR\ refuted (attributed to activity, an alias or the sampling, or not
recovered with better data) and \nmBenchD\ disputed, which we exclude from the metrics.
Of these signals, \nmBenchNotHarps\ were not discovered with HARPS data. The labels, periods
and references were compiled with web searches by an AI agent and each was checked
against the publisher's record; we did not choose signals by how the ladder treats
them. For every signal we remove the other accepted planets of its system with
Keplerian fits and evaluate the ladder at the claimed period.

Three caveats apply. Several refuted signals were refuted with the same kinds of
tests the ladder uses, so agreement is partly circular. The refuted signals were
claimed from subsets of the data now in the archive, so the benchmark tests the
ladder on the full archive, not on the data the discoverers had. And the numbers are
small, so we give exact binomial intervals.

\subsection{Significance}

Whether a claimed signal is significant depends on the question. A blind search of
the full archive needs the global FAP over the search band: by that measure
\nmRefNotDet\ of the \nmBenchR\ refuted signals are not significant (FAP $> 1\%$), while
\nmUpDet\ of the \nmUpN\ upheld signals are (Fig.~\ref{fig:bench}). The \nmUpNotDet\
upheld signals that are not have measured $K/\sigma_K$ of
\nmUpNotDetSnrLo--\nmUpNotDetSnrHi\ in these data, all were confirmed with other
instruments or by transits, and all lie below the 50\% completeness point of
Fig.~\ref{fig:completeness}, although that point is in expected $K/\sigma_K$, about which
the measured value scatters by roughly one. For a period known in advance,
the single-frequency probability is the right measure; below 1\% by that measure,
\nmRefNotDetSingle\ refuted and \nmUpNotDetSingle\ upheld signals are not significant.
This probability is taken at the best point within $\pm0.5/T$ and from the likelihood
ratio, which Sect.~\ref{sec:fap} found liberal, so it overstates significance somewhat;
Table~\ref{tab:window} gives the counts at the claimed period itself. Even so, most
refuted signals leave some power at their period in the full archive; what they lack is
the strength to stand out in a blind search.

\subsection{The rule and the classifiers}

Among significant signals, the a-priori rule accepts \nmBRuleRefAcc\ of the
\nmBRuleRefN\ refuted ones (95\% interval \nmBRuleFPRLo--\nmBRuleFPRHi\%) and rejects
\nmBRuleUpRej\ of the \nmBRuleUpN\ upheld ones (\nmBRuleFNRLo--\nmBRuleFNRHi\%):
\nmBRuleUIRej\ of \nmBRuleUIN\ upheld by independent evidence and \nmBRuleUHRej\ of
\nmBRuleUHN\ accepted from HARPS alone. Table~\ref{tab:refsig} lists the features of the
significant refuted signals and of any significant upheld signal the rule rejects. All
\nmTopUp\ significant upheld signals are the highest peak in the band, against
\nmTopRef\ of the \nmBRuleRefN\ significant refuted ones, a difference the classifiers,
trained on highest peaks, do not see.
GJ\,581\,d, Barnard's star b at 233\,d and HD\,41248\,c fail on indicator power (the first
was attributed to activity by \citealt{Robertson2014}, the second to a yearly alias of
rotation by \citealt{Lubin2021}, the third to rotation by \citealt{Santos2014}, whose
rotation period our proxy recovers); Kapteyn c fails on phase and has H$\alpha$ power as
well; AD\,Leo\,b fails the period rung, because a stronger peak of the same activity signal
lies a few resolution elements away, and has chromatic-index power, as expected for the
spot-induced signal of \citet{Carleo2020}. The rule accepts HD\,26965\,b, for which the
indicators show neither power above the threshold at the period nor a significant rotation
proxy. It rejects two upheld signals, HD\,192310\,b and GJ\,163\,c, on indicator power.
Several margins are thin (Table~\ref{tab:refsig}). Adding the rotation-alias test of the first benchmark pass (\code{harm\_alias\_dist}) to
the rule \nmBRuleAliasEffect.

The classifiers rank the significant signals with an AUC of \nmBAucTree\ (trees; 95\%
interval \nmBAucTreeLo--\nmBAucTreeHi\ from a bootstrap within each class) and
\nmBAucLog\ (logistic regression; \nmBAucLogLo--\nmBAucLogHi), misordering \nmBMisTree\
and \nmBMisLog\ of the \nmBPairs\ upheld--refuted pairs. At a score of 0.5 the trees
accept \nmBTreeRefAcc\ significant refuted signal (GJ\,581\,d) and reject
\nmBTreeUpRej\ upheld ones (HD\,15337\,c and HD\,192310\,b); the logistic regression accepts
\nmBLogRefAcc\ (GJ\,581\,d) and rejects \nmBLogUpRej\ (HD\,15337\,c, the weakest upheld signal
that is significant).
GJ\,581\,d has indicator power above the rule's threshold but a strong, stationary,
alias-free peak, and both models weigh the latter more. Without the \nmBEasyN\ upheld
signals with $K/\sigma_K > 20$ the trees' AUC is \nmBAucTreeNoEasy; for signals upheld by
independent evidence only, \nmBAucTreeUI. The model without indicator features reaches
\nmBAucNoInd\ (\nmBAucNoIndLo--\nmBAucNoIndHi). Among the same significant signals,
single features reach AUCs of \nmBSingleGrowth\ (growth of $\dchi$), \nmBSingleAlias\
(alias margin), \nmBSingleBlockRms\ (block phase rms), \nmBSingleInd\ (indicator power)
and \nmBSingleFap\ (FAP). On signals that are not significant, the classifiers, trained
only on significant peaks, are outside their training domain and their scores should
not be used.

Only \nmRefOutsideProt\ refuted host, AD\,Leo, has a published rotation period outside the
simulated 3--150\,d; at 2.23\,d it is also below the 2.5-d lower limit of the rotation
proxy, so no rotation test is possible for it.

The rotation proxy behaves differently on real stars. It is significant for
\nmProxyHostsN\ of the \nmNHosts\ benchmark hosts and for \nmProxyUpSigN\ of the
\nmBRuleUpN\ significant upheld signals, but for only \nmProxyPlanetPct\% of simulated
planets: permutation destroys the star's own coherent indicator variations, rotation and
cycles alike, so simulated planets sit on indicators without activity and real ones do
not. Where a host has a significant proxy and a published rotation period, the two agree
within 20\% for \nmProxyAgreeN\ of \nmProxyCmpN. The rotation features of real candidates
therefore lie partly outside the simulated range. A model trained
without the rotation-proxy features reaches an AUC of \nmBAucNoProxy\
(\nmBAucNoProxyLo--\nmBAucNoProxyHi) on the significant benchmark signals.

\begin{table*}
\caption{Significant refuted signals, and significant upheld signals the a-priori rule
rejects: the features behind each verdict.}
\label{tab:refsig}
\centering
\footnotesize
\begin{tabular}{lcrrlrrlccc}
\hline\hline
Signal & Verdict & $\log_{10}$FAP & Alias margin & Indicator $\dchi$ & $P_{\rm proxy}$ (d) & \code{harm\_dist} & Fails at & Trees & Logistic & No ind. \\
\hline
\inputrows{tab_refsig_rows}
\hline
\end{tabular}
\tablefoot{Alias margin: $\dchi(f_0)$ minus the strongest alias inside the band. Fails at:
the first rung of the a-priori rule the signal fails (period, phase; ind.: indicators;
rot.: rotation; --: passes).
Indicator $\dchi$: the largest near $f_0$, with the indicator. $P_{\rm proxy}$: the
rotation proxy, when significant. The last three columns are classifier scores.}
\end{table*}

\begin{figure}
\centering
\includegraphics[width=\columnwidth]{fig_benchmark.pdf}
\caption{The benchmark: score from the boosted trees for signals upheld by
independent evidence (blue), accepted from HARPS data (green), refuted (red) and
disputed (grey). Open symbols are not significant at the claimed period in a blind
search of the full archive; crosses mark significant signals rejected by the a-priori
rule.}
\label{fig:bench}
\end{figure}

\subsection{Changes after the first pass}
\label{sec:changes}

The rule was frozen before the first benchmark pass, and its thresholds have not
changed. The pipeline has, four times. After the first pass exposed failures we lowered
the S/N cut from 20 to 10 (it removed \nmProxLossTwenty\% of the spectra of Proxima
Centauri), replaced two-term Fourier prewhitening of known companions by Keplerian
orbits (the Fourier terms failed for $\pi$\,Men\,b with $e = 0.64$), added the quadratic
trend and the 15-$\sigma$ clip (for the binaries $\alpha$\,Cen\,B and GJ\,667\,C), and
added \code{harm\_alias\_dist}. An independent review of the first draft then found
four errors, which version 0.2 corrected: nights began at 0\,h\,UT instead of
local noon; the period rung compared candidates with aliases outside the search band
and could refine onto the candidate's own peak, so that it rejected planets on their
one-day mirrors; the jitter was held at its planet-free value; and the block tests had
the wrong degrees of freedom. At the same time we widened the window around a claimed
period from $\pm0.25/T$ to $\pm0.5/T$, after the review showed a refuted signal refined to
the edge of the narrower window, and changed the simulations (Sect.~\ref{sec:sims}). A
second review found that the effective window had grown to $\pm0.6/T$, now capped at
$\pm0.5/T$, and that the rotation proxy fired on slow variations of real indicators; the
indicators are now detrended like the velocities. A third review found that this
detrending, done by least squares, was tilted by indicator outliers in the simulations,
and that many real proxies sat at half a year; outliers are now clipped before the fit and
the seasonal aliases are excluded from the proxy search (version \nmrvvetVersion).
Table~\ref{tab:window} shows how the benchmark results depend on the window: with the
first-pass window the rule rejects \nmWinQRefRej\ of \nmWinQRefSig\ significant refuted
signals and \nmWinQUpRej\ of \nmWinQUpSig\ upheld ones; at the claimed period itself,
\nmWinZRefRej\ of \nmWinZRefSig\ and \nmWinZUpRej\ of \nmWinZUpSig.
Of the benchmark signals, \nmNEdge\ reach the edge of the adopted window (marked in
Table~\ref{tab:bench}).

\begin{table*}
\caption{Benchmark results for three windows around the claimed period.}
\label{tab:window}
\centering
\footnotesize
\begin{tabular}{lcccccccc}
\hline\hline
 & \multicolumn{2}{c}{Refuted, significant} & \multicolumn{2}{c}{Upheld, significant} & \multicolumn{2}{c}{AUC (significant)} & \multicolumn{2}{c}{Not significant, single $f$} \\
Window & $N$ & rule rejects & $N$ & rule rejects & trees & logistic & refuted & upheld \\
\hline
\inputrows{tab_window_rows}
\hline
\end{tabular}
\tablefoot{Significant: global Baluev FAP $< 1\%$. The last two columns count signals whose
single-frequency probability is above 1\% (of \nmBenchR\ refuted and \nmUpN\ upheld).}
\end{table*} Table~\ref{tab:changes} lists the \nmBChanged\ benchmark signals
whose outcome differs between the first pass and this version. The frozen-pipeline
results are released with the code.

\subsection{A test forward in time}
\label{sec:timesplit}

The benchmark asks how the ladder treats signals whose fate was already known. A stricter
test asks whether it would have pointed the right way before the fate was known.
HARPS-RVBank holds spectra taken before 2022 (the last on 2021 December 29), so we ran the
frozen ladder blind on every star with at least 20 nights over 100\,d (\nmTsStars\ stars;
the highest peak per star, 1.2--500\,d) and built the list it would have produced at the
start of 2022: significant peaks that pass the a-priori rule, have $K < 100\ms$ and lie on
stars with no planet known before 2022. We count a planet as known before 2022 if the NASA
Exoplanet Archive dates its discovery paper before 2022 (queried on 2026 October 4) or
exoplanet.eu dates its discovery before 2022 (queried on 2026 October 6). The second
catalogue also counts earlier announcements, such as planets announced in 2011
\citep{Mayor2011} that the NASA archive lists only from a later paper or not at all. That
list has \nmTsN\ signals. Since then, the NASA archive has listed \nmTsHits\ of them as
planets published at the same period ($|1/P - 1/P_{\rm pub}| < 1.5/T$).

Ranked by the boosted trees, \nmTsTreeTen\ of the top 10 and \nmTsTreeTwenty\ of the top 20
were later published, against \nmTsExpTen\ and \nmTsExpTwenty\ expected from the base rate
of \nmTsBase\% (hypergeometric $p = \nmTsPTen$ and $\nmTsPTwenty$;
Fig.~\ref{fig:timesplit}). Ranked by the false-alarm probability alone the counts are
\nmTsFapTen\ and \nmTsFapTwenty, and by the logistic regression \nmTsLogTen\ and
\nmTsLogTwenty. Over the whole list the area under the ROC curve is \nmTsAucTree\ (95\%
interval \nmTsAucTreeLo--\nmTsAucTreeHi) for the trees, \nmTsAucLog\
(\nmTsAucLogLo--\nmTsAucLogHi) for the logistic regression, \nmTsAucFap\
(\nmTsAucFapLo--\nmTsAucFapHi) for the false-alarm probability and \nmTsAucSnr\ for
$K/\sigma_K$. The difference between the trees and the false-alarm probability is
\nmTsDauc\ (paired bootstrap, 95\% interval \nmTsDaucLo\ to \nmTsDaucHi). The later planets gather at the top of the list, but with
\nmTsHits\ positives this test does not show that the classifiers rank better than
significance alone.

The concentration at the top survives the other choices we tried. Keeping stars that had
other planets before 2022 (\nmTsLooseN\ signals, the same \nmTsLooseHits\ later planets)
leaves the trees at \nmTsLooseTreeTen\ of the top 10. Dating planets by the NASA archive
alone keeps \nmTsEuOnlyN\ more stars, among them
HD\,134606 and HD\,137496, whose planets were announced in 2011 and 2021, and the hosts of
two TESS planets that exoplanet.eu dates to 2019; it gives \nmTsNeaHits\ later planets among
\nmTsNeaN\ signals, \nmTsNeaTreeTen\ of them in the trees' top 10. A window of $1/T$ instead
of $1.5/T$ changes nothing. On short baselines the window is wide, though: \nmTsWideN\ of
the later planets differ from the archival period by \nmTsWideMin--\nmTsWideMax\%. Requiring
the periods to agree within 5\% leaves \nmTsFiveHits\ later planets, \nmTsFiveTreeTen\ of them
in the trees' top 10 (\nmTsFiveExpTen\ expected, $p = \nmTsFivePTen$) and \nmTsFiveFapTen\ in
the top 10 by the false-alarm probability.

Two limits apply. Of the \nmTsHits\ later planets, \nmTsToi\ orbit stars with TESS
candidates, often observed with HARPS because of the candidate, so the list is enriched in
stars someone was already following; \nmTsToiTopTen\ of the \nmTsTreeTen\ in the trees' top
10 are such stars, and without them \nmTsNonToi\ later planets remain, too few to test. And
``not yet published'' is not ``false'': the share of later planets is a lower bound that
may rise as the other signals are followed up.

\begin{figure}
\centering
\includegraphics[width=\columnwidth]{fig_timesplit.pdf}
\caption{The \nmTsN\ signals the ladder would have listed in January 2022 on stars with no
planet known before 2022, ranked by each score: the cumulative number published as planets at the same
period since then (\nmTsHits\ in all).}
\label{fig:timesplit}
\end{figure}

The open signals also make predictions that new data can test. We froze them, with
circular ephemerides, in a list (the discovery ledger) and tested the two of its 30
highest-ranked signals that had at least ten public spectra taken after 2022 in the ESO
archive, with the decision rule written down before their velocities were read. For
GJ\,902 (\nmTaGjP\,d, predicted $K = \nmTaGjK \pm \nmTaGjsK\ms$), ESPRESSO, HARPS and NIRPS
spectra on \nmTaGjNights\ nights (\nmTaGjVel\ nightly velocities, one per instrument and night)
give $\nmTaGjA \pm \nmTaGjsA\ms$ at the predicted phase. One ESPRESSO product in the search
cone belongs to another star ($-106$ against $+70$\,km\,s$^{-1}$) and was left out, a
decision taken after the headers were read. The ladder and the rule passed this signal
(trees score \nmTaGjTree), and the new data exclude it. For HD\,58489 (\nmTaHdP\,d,
$\nmTaHdK \pm \nmTaHdsK\ms$), \nmTaHdNights\ HARPS nights give $\nmTaHdA \pm \nmTaHdsA\ms$ at the
predicted phase, about $\nmTaHdSigma\sigma$ once the statistic is calibrated on noise-only
simulations. That result gives the two cross-correlation masks of these spectra separate
zero points, a choice made after the velocities were displayed but before any fit. With
the single zero point that was registered, the amplitude is $\nmTaHdLitA \pm \nmTaHdLitsA\ms$
and the test is undecided, which is the registered verdict. Even with separate zero points,
a sinusoid at a random phase does as well in \nmTaHdPhase\% of trials, and \nmTaHdPer\% of
periods between 2 and 150\,d fit the new velocities better. One prediction failed and the
other is undecided. The ledger, the rule, the data and the results are released with the
code; git commit times record the order in which they were written, but they are not
independent timestamps.

%%%%%%%%%%%%%%%%%%%%%%%%%%%%%%%%%%%%%%%%%%%%%%%%%%%%%%%%%%%%%%%%%%%%%%%%%%%%%%%
\section{\texorpdfstring{A worked example: HD\,297396\,\ensuremath{\mathrm{b}}}{A worked example: HD 297396 b}}
\label{sec:example}

HD\,297396\,b is a 4.27-d candidate found in a periodogram survey of HARPS-RVBank \citep{Fraser2026}. Through the generic pipeline, with \nmHDn\ nightly epochs over
\nmHDbase\,yr (the discovery analysis, with stricter cuts, keeps 104), its peak has a
Baluev FAP of $10^{\nmHDfap}$ and $K = \nmHDK \pm \nmHDsK\ms$ ($K/\sigma_K = \nmHDsnr$).
It beats its aliases by $\dchi = \nmHDalias$, holds phase across halves
($\nmHDhalf\sigma$) and blocks ($p = \nmHDblockp$), grows steadily with the number of
epochs ($\rho = \nmHDrho$), gains $\dchi = \nmHDapod$ from an envelope, and no indicator
has power near its period beyond $\dchi = \nmHDind$. The indicators give no significant
rotation proxy, so the rotation rung is not applied.
It \nmHDrule\ the a-priori rule, whose thresholds were partly set during its own
analysis, and the classifiers give it scores of \nmHDpTree\ (trees), \nmHDpLog\
(logistic) and \nmHDpNoInd\ (without indicators). These are not the candidate's adopted
values. The discovery analysis \citep{Fraser2026} gives each of four groups of
programmes its own offset and jitter and finds $K = 5.5 \pm 0.8\ms$. Its false-alarm
probability, calibrated by simulation, is $1.4\times10^{-3}$ with all epochs; without one
discrepant night no simulation exceeds the peak, and the analytic bound is $7\times10^{-6}$.
Most of the gap comes from the jitter. The generic statistic refits the jitter scale with
the signal in the model; with the jitters held at their planet-free values, the same
pipeline gives $10^{\nmHDfapHz}$ with all epochs and $10^{\nmHDfapHzNoNight}$ without the
discrepant night (Sect.~\ref{sec:fap} found the first treatment liberal and the second
conservative on permuted real noise). We do not break down the remaining factor of
\nmHDratioLo--\nmHDratioHi: the two analyses also differ in their noise models (two eras
against four groups of programmes), their epochs and their frequency bands.

The example also shows what the ladder cannot do. The star's rotation period is not
measured; the indicators are weak for this star; and activity calibrations place its
rotation near 39\,d, close to nine times the orbital period, a harmonic that our
simulated activity rarely puts power into \citep{Fraser2026}. A high score says
that the signal behaves like the simulated planets and unlike the simulated
impostors; it does not replace a second instrument.
